

The digital appendix contains the thesis documents, the prototypical
implementation, and the experimental materials supporting this work.
It provides access to the implementation described in Chapter~\ref{chap:implementation},
the evidence supporting Chapter~\ref{chap:evaluation}, and an inspectable copy of the
constructed Neo4j graph.

The materials follow the organization of the project repository:

\begin{verbatim}
defi-verification/
|-- thesis/
|   |-- DEFI_VERIFICATION.pdf
|   |-- source/
|   |-- figures/
|-- source/
|-- evaluation/
|   |-- 400_prompts/
|   |-- 600_prompts/
|   |-- datasets/
|   |-- graph/
|   |   |-- submission/
|   |       |-- database/
|   |       |-- evidence/
|   |       |-- figures/
|   |       |-- queries/
|   |       |-- README.md
|   |-- sources/
|   |-- tests/
|   |-- wallet/
|-- README.md
|-- requirements.txt
|-- requirements-dev.txt
|-- pytest.ini
|-- .env.example
\end{verbatim}

The directory listing shows the principal components. Additional
scripts, documentation, and integrity manifests are provided within
the corresponding directories.

\section{Thesis Documents}

The \texttt{thesis/} directory contains the final thesis in PDF format,
corresponding to the submitted printed version. The
\texttt{thesis/source/} directory contains the editable thesis project,
bibliography, and supporting files required to generate the document.
Editable diagram sources and figure-generation scripts are included
in \texttt{thesis/figures/}.

\section{Prototypical Implementation}

The \texttt{source/} directory contains the implementation of the
graph-based intent-verification system. It includes natural-language
intent extraction, normalization and validation, authoritative
knowledge acquisition, graph synchronization, relationship
verification, and the wallet execution adapter. The browser
interface and application tests are also included.

The root \texttt{README.md} documents software requirements,
installation, environment configuration, application startup, and
testing. Python dependencies are specified in
\texttt{requirements.txt} and \texttt{requirements-dev.txt}.
Node.js dependencies and their lock file are provided in
\texttt{source/wallet/}. The \texttt{.env.example} file describes
the required configuration without including secret credentials.

\section{Experimental Materials and Results}

The \texttt{evaluation/} directory contains evaluation runners,
datasets, reference information, raw observations, and analysis
outputs supporting Chapter~5.

\begin{description}
    \item[\texttt{400\_prompts/}]
    Materials from the exploratory 400-prompt study, including
    frozen inputs, recorded interpretations, verification outcomes,
    execution journals, and analysis outputs.

    \item[\texttt{600\_prompts/}]
    Materials from the balanced 600-prompt study, including
    the frozen dataset, execution order, recorded observations,
    error reviews, metrics, tables, and figures.

    \item[\texttt{datasets/}]
    Additional evaluation datasets and controlled test cases.

    \item[\texttt{graph/}]
    Evidence from graph-verification, controlled
    relationship-mutation, knowledge-maintenance, and traceability
    studies. The \texttt{submission/} subdirectory additionally
    provides the database export and inspection materials described
    below.

    \item[\texttt{wallet/}]
    Results from the supplementary controlled wallet evaluation.
    These observations concern policy and execution-boundary
    behaviour; simulated outcomes do not establish successful
    on-chain execution.

    \item[\texttt{sources/}]
    Source provenance, historical implementation snapshots,
    environment records, and supporting reproduction materials.

    \item[\texttt{tests/}]
    Tests for evaluation runners, datasets, metrics, and
    evidence-integrity checks.
\end{description}

Raw observations include recorded model interpretations, provider
attempts and errors, verification decisions, graph events, and
timings. Derived metrics and analysis outputs are retained alongside
the corresponding experimental records.

\section{Constructed Neo4j Graph}

The directory \texttt{evaluation/graph/submission/} contains
materials for inspecting the graph used by the prototype.
The submission snapshot was exported on 11 November 2026 using
Neo4j 2026.05.0 Enterprise.

\begin{description}
    \item[\texttt{database/}]
    Contains the native \texttt{neo4j.dump} database export and
    metadata recording its version, checksum, and validation.
    The export contains the \texttt{neo4j} user database;
    the system database and other local databases are excluded.

    \item[\texttt{queries/}]
    Contains read-only Cypher queries for displaying the protocol
    and deployment overview, inspecting an Aave V3 USDC supply
    subgraph, counting nodes and relationships, and examining
    knowledge provenance and synchronization timestamps.

    \item[\texttt{evidence/}]
    Contains saved query results, database counts, selected node
    and relationship properties, and a recorded response from the
    relationship verifier.

    \item[\texttt{figures/}]
    Contains PNG and SVG diagrams generated from the saved query
    results. These figures present selected stored relationships;
    they are not screenshots of the Neo4j interface or
    visualizations of the complete database.
\end{description}

The inspected database contains 3,736 nodes and 20,729 relationships,
including presentation and fixture metadata. The focused Aave supply
view contains eight nodes and nine relationships connecting the
network, protocol deployment, authorized contract and function,
supported action, reserve, and underlying token.

The recorded verifier response is
\texttt{INSUFFICIENT\_EVIDENCE}, with rule
\texttt{KNOWLEDGE.STALE}, because the stored knowledge exceeded the
configured freshness limit at inspection time. The structural
relationships remain visible, but their presence alone does not
establish current approval. The freshness policy and synchronization
timestamps were not modified to obtain a different result.

The submission snapshot is preserved separately from the historical
experimental evidence. It is not claimed to be the exact database
state used for the earlier evaluations. The export was successfully
loaded into an isolated data directory and passed an offline
consistency check. Restoration and inspection instructions are
provided in the accompanying \texttt{README.md}.

\section{Inspection and Reproduction}

The root \texttt{README.md} explains how to configure and run the
prototype. The file \texttt{evaluation/README.md} describes how to
inspect the retained experimental evidence, verify its integrity,
and execute the evaluation procedures.

Saved reports, observations, and graph figures can be inspected
without rerunning the experiments. Live evaluation requires the
documented external services, configuration, and graph prerequisites.
Repeated executions may produce different observations because
model responses, external knowledge sources, and blockchain state
can change.

Historical source snapshots and results retain their original
contents. Paths within these artifacts may therefore reflect the
directory structure used when the experiments were conducted.
New experiments should write their outputs to separate locations
rather than overwrite the retained evidence.